%=============================================================================
% CHAPTER 2 — COMPLETE PEDAGOGICAL ARCHITECTURE
% Statistical Distributions for Financial Markets
%
% Three-Layer Design:
%   Layer 1 — Year 3 (Undergraduate Core)
%   Layer 2 — Master Extension
%   Layer 3 — PhD / Research Extension
%
% Author: Daniel Traian PELE
% Course: Statistics of Financial Markets
%=============================================================================

\documentclass[11pt, a4paper]{article}
\usepackage[utf8]{inputenc}
\usepackage[T1]{fontenc}
\usepackage{amsmath, amssymb, amsthm}
\usepackage{mathtools}
\usepackage{bm}
\usepackage{booktabs}
\usepackage{array}
\usepackage{geometry}
\usepackage{enumitem}
\usepackage{hyperref}
\usepackage{xcolor}

\geometry{margin=2.5cm}
\definecolor{MainBlue}{RGB}{26, 58, 110}
\definecolor{Crimson}{RGB}{220, 53, 69}
\definecolor{Forest}{RGB}{46, 125, 50}
\definecolor{Purple}{RGB}{142, 68, 173}
\hypersetup{colorlinks=true, linkcolor=MainBlue, urlcolor=MainBlue}

\newcommand{\E}{\mathbb{E}}
\newcommand{\Var}{\text{Var}}
\newcommand{\Cov}{\text{Cov}}
\newcommand{\Corr}{\text{Corr}}
\newcommand{\R}{\mathbb{R}}
\newcommand{\layerY}[1]{\textcolor{MainBlue}{\textbf{[Y3]}} #1}
\newcommand{\layerM}[1]{\textcolor{Forest}{\textbf{[MSc]}} #1}
\newcommand{\layerP}[1]{\textcolor{Purple}{\textbf{[PhD]}} #1}

\theoremstyle{definition}
\newtheorem{defn}{Definition}
\newtheorem{thm}{Theorem}
\newtheorem{prop}{Proposition}

\title{\textbf{Chapter 2: Statistical Distributions for Financial Markets}\\[6pt]
\large Complete Pedagogical Architecture\\[4pt]
\normalsize Statistics of Financial Markets}
\author{Daniel Traian PELE}
\date{}

\begin{document}
\maketitle
\tableofcontents
\newpage

%=============================================================================
%=============================================================================
\part{CHAPTER ARCHITECTURE}
%=============================================================================
%=============================================================================

%#############################################################################
\section{Hierarchical Chapter Outline}
%#############################################################################

\begin{enumerate}[label=\textbf{2.\arabic*}, leftmargin=2cm]

\item \textbf{Introduction — Why Distributions Matter in Finance}
  \begin{itemize}[nosep]
    \item[\layerY{}] Risk measurement motivation, tail risk intuition, model risk
    \item[\layerM{}] Distribution choice and derivative pricing
    \item[\layerP{}] Model uncertainty quantification
  \end{itemize}

\item \textbf{From Prices to Log-Returns}
  \begin{itemize}[nosep]
    \item[\layerY{}] Definition, computation, statistical properties
    \item[\layerM{}] Continuous-time limit, Itô calculus connection
    \item[\layerP{}] Semimartingale representation
  \end{itemize}

\item \textbf{The Normal Distribution — The Classical Model}
  \begin{itemize}[nosep]
    \item[\layerY{}] PDF, moments, CLT, quantiles (preparation for VaR in Section 2.10)
    \item[\layerM{}] Maximum entropy characterization, normality tests
    \item[\layerP{}] Berry--Esseen bounds, rate of convergence
  \end{itemize}

\item \textbf{Stylized Facts of Financial Returns}
  \begin{itemize}[nosep]
    \item[\layerY{}] Fat tails, skewness, volatility clustering
    \item[\layerM{}] Aggregational Gaussianity, leverage effect, long memory
    \item[\layerP{}] Cont (2001) taxonomy, fractal structure, multiscaling
  \end{itemize}

\item \textbf{The Student-$t$ Distribution}
  \begin{itemize}[nosep]
    \item[\layerY{}] PDF, degrees of freedom, excess kurtosis, MLE fitting
    \item[\layerM{}] Location-scale form, AIC/BIC comparison, GARCH-$t$
    \item[\layerP{}] Mixture representation, connection to F-distribution, Bayesian interpretation
  \end{itemize}

\item \textbf{Asymmetric Distributions}
  \begin{itemize}[nosep]
    \item[\layerY{}] Skewness effect on downside risk, basic intuition
    \item[\layerM{}] Hansen's skewed-$t$, GED, MLE estimation
    \item[\layerP{}] Generalized Hyperbolic family, NIG, Barndorff-Nielsen (1977)
  \end{itemize}

\item \textbf{Stable ($\alpha$-Stable) Distributions} \textbf{[MANDATORY]}
  \begin{itemize}[nosep]
    \item[\layerY{}] Intuition, stability property, tail index $\alpha$, why variance may not exist
    \item[\layerM{}] Characteristic function, GCLT, estimation (McCulloch, MLE), financial interpretation
    \item[\layerP{}] Regular variation, domain of attraction, Lévy processes, subordinators, tempered stable
  \end{itemize}

\item \textbf{Extreme Value Theory (EVT)}
  \begin{itemize}[nosep]
    \item[\layerY{}] Tail behavior intuition, GPD concept (descriptive), extreme events in finance
    \item[\layerM{}] Fisher--Tippett--Gnedenko theorem, POT method, threshold selection, EVT-based VaR
    \item[\layerP{}] Hill estimator, mean excess function, second-order conditions, bias--variance trade-off
  \end{itemize}

\item \textbf{Distribution Comparison and Model Selection}
  \begin{itemize}[nosep]
    \item[\layerY{}] Visual comparison, QQ-plots, basic KS test
    \item[\layerM{}] AIC/BIC, Anderson--Darling, likelihood ratio tests
    \item[\layerP{}] Vuong test, cross-validation, scoring rules, Diebold--Mariano
  \end{itemize}

\item \textbf{Risk Management Implications}
  \begin{itemize}[nosep]
    \item[\layerY{}] \textbf{Define VaR from scratch}: definition, interpretation, graphical intuition, VaR under normality, numerical example, comparison Normal vs.\ Student-$t$
    \item[\layerM{}] \textbf{Define ES from scratch}, coherence axioms, Kupiec test, backtesting, Basel framework
    \item[\layerP{}] Spectral risk measures, elicitability, regulatory capital optimization
  \end{itemize}

\item \textbf{Modern Extensions}
  \begin{itemize}[nosep]
    \item[\layerY{}] GARCH intuition (conceptual only)
    \item[\layerM{}] GARCH-$t$, Gaussian mixtures, regime switching
    \item[\layerP{}] Stochastic volatility with jumps, Lévy-driven models, realized measures
  \end{itemize}

\end{enumerate}

%#############################################################################
\section{Slide Count Allocation}
%#############################################################################

\begin{center}
\renewcommand{\arraystretch}{1.3}
\begin{tabular}{l c c c c}
\toprule
\textbf{Section} & \textbf{Y3 Slides} & \textbf{MSc Slides} & \textbf{PhD Slides} & \textbf{Total} \\
\midrule
2.1 Introduction                  & 3  & 1  & 1  & 5  \\
2.2 Prices to Log-Returns         & 2  & 1  & 1  & 4  \\
2.3 Normal Distribution           & 4  & 2  & 1  & 7  \\
2.4 Stylized Facts                & 4  & 2  & 2  & 8  \\
2.5 Student-$t$ Distribution      & 4  & 2  & 1  & 7  \\
2.6 Asymmetric Distributions      & 2  & 3  & 2  & 7  \\
2.7 Stable Distributions          & 4  & 4  & 4  & 12 \\
2.8 Extreme Value Theory          & 3  & 4  & 3  & 10 \\
2.9 Distribution Comparison       & 3  & 2  & 2  & 7  \\
2.10 Risk Management              & 5  & 3  & 2  & 10 \\
2.11 Modern Extensions            & 1  & 3  & 3  & 7  \\
Conclusions + Homework            & 2  & 1  & 1  & 4  \\
\midrule
\textbf{Total}                    & \textbf{37} & \textbf{28} & \textbf{23} & \textbf{88} \\
\bottomrule
\end{tabular}
\end{center}

%=============================================================================
%=============================================================================
\part{LAYERED CONTENT STRUCTURE}
%=============================================================================
%=============================================================================

%#############################################################################
\section{Section 2.1: Introduction — Why Distributions Matter}
%#############################################################################

\subsection{Year 3 Core}
\begin{itemize}
  \item \textbf{Motivating question}: Why does it matter which probability distribution we use to model financial returns?
  \item If we assume the wrong distribution, we underestimate the probability of extreme losses
  \item \textbf{Example}: October 19, 1987 — the S\&P 500 fell 20.5\% in a single day
  \begin{itemize}
    \item Under the Normal assumption, this would be a $\sim$25$\sigma$ event with probability $\approx 10^{-135}$
    \item Yet it happened — suggesting the Normal distribution is inadequate for tail events
  \end{itemize}
  \item Banks and regulators need to quantify ``how bad can it get?'' — this requires accurate distributional models
  \item \textbf{Preview}: In Section 2.10 we will formally define \textbf{Value-at-Risk (VaR)}, a key risk measure that depends entirely on the distributional assumption
  \item The distribution choice problem: Normal vs.\ Student-$t$ vs.\ Stable vs.\ EVT
\end{itemize}

\subsection{Master Extension}
\begin{itemize}
  \item Impact on derivative pricing: Black--Scholes assumes log-normality; deviations produce the volatility smile
  \item Portfolio optimization under non-Gaussian returns: mean-variance is suboptimal
  \item Regulatory context: Basel II $\to$ Basel III/IV framework (details in Section 2.10)
\end{itemize}

\subsection{PhD Extension}
\begin{itemize}
  \item Model uncertainty and ambiguity aversion (Ellsberg paradox in portfolio choice)
  \item Robust risk measures under distributional uncertainty
  \item Minimax approaches: worst-case distributional analysis
\end{itemize}

%#############################################################################
\section{Section 2.2: From Prices to Log-Returns}
%#############################################################################

\subsection{Year 3 Core}

\begin{defn}[Log-Return]
The \textbf{log-return} (continuously compounded return) at time $t$ is:
\[
  r_t = \ln\!\left(\frac{P_t}{P_{t-1}}\right) = \ln P_t - \ln P_{t-1}
\]
where $P_t$ is the asset price at time $t$.
\end{defn}

\begin{itemize}
  \item \textbf{Additivity}: Multi-period returns are sums: $r_{t:t+k} = \sum_{i=1}^{k} r_{t+i}$
  \item \textbf{Approximation}: For small returns, $r_t \approx (P_t - P_{t-1})/P_{t-1}$
  \item \textbf{Statistical motivation}: Additive structure enables standard probability theory
\end{itemize}

\subsection{Master Extension}
\begin{itemize}
  \item Continuous-time limit: $dS/S = \mu\,dt + \sigma\,dW$ $\Rightarrow$ $d\ln S = (\mu - \sigma^2/2)\,dt + \sigma\,dW$
  \item Jensen's inequality: $\E[\ln(P_T/P_0)] \neq \ln(\E[P_T/P_0])$
  \item Log-return vs.\ simple return: implications for portfolio aggregation
\end{itemize}

\subsection{PhD Extension}
\begin{itemize}
  \item Semimartingale decomposition and Itô's lemma
  \item Quadratic variation and realized volatility: $[X,X]_t = \lim \sum (X_{t_{i+1}} - X_{t_i})^2$
  \item Non-continuous price processes: jumps and Lévy--Itô decomposition
\end{itemize}

%#############################################################################
\section{Section 2.3: The Normal Distribution}
%#############################################################################

\subsection{Year 3 Core}

\begin{defn}[Normal Distribution]
$X \sim N(\mu, \sigma^2)$ if:
\[
  f(x) = \frac{1}{\sigma\sqrt{2\pi}}\exp\!\left(-\frac{(x-\mu)^2}{2\sigma^2}\right), \quad x \in \R.
\]
\end{defn}

\noindent\textbf{Key properties}:
\begin{itemize}
  \item $\E[X] = \mu$, \quad $\Var(X) = \sigma^2$
  \item Skewness $= 0$, \quad Kurtosis $= 3$ (excess kurtosis $= 0$)
  \item 68--95--99.7 rule
  \item \textbf{Quantiles}: The $\alpha$-quantile of $N(\mu, \sigma^2)$ is $q_\alpha = \mu + \sigma\,\Phi^{-1}(\alpha)$; this will be used in Section 2.10 to define Value-at-Risk
  \item \textbf{CLT}: $\frac{\bar{X}_n - \mu}{\sigma/\sqrt{n}} \xrightarrow{d} N(0,1)$
  \item \textbf{Stability}: $X + Y \sim N(\mu_1+\mu_2, \sigma_1^2+\sigma_2^2)$ for independent normals
  \item \textbf{Limitation}: Tails decay as $e^{-x^2/2}$ — too thin for financial data
\end{itemize}

\subsection{Master Extension}
\begin{itemize}
  \item Maximum entropy characterization: among all distributions with given $(\mu, \sigma^2)$, the normal maximizes Shannon entropy
  \item Normality tests: Jarque--Bera, Shapiro--Wilk, Anderson--Darling, Kolmogorov--Smirnov
  \[
    JB = \frac{n}{6}\left(S^2 + \frac{(K-3)^2}{4}\right) \sim \chi^2(2)
  \]
  \item Moment generating function: $M_X(t) = \exp(\mu t + \sigma^2 t^2 / 2)$
  \item Multivariate normal and Mahalanobis distance
\end{itemize}

\subsection{PhD Extension}
\begin{itemize}
  \item Berry--Esseen theorem: $\sup_x |F_n(x) - \Phi(x)| \leq \frac{C\,\rho}{\sigma^3\sqrt{n}}$, where $\rho = \E[|X-\mu|^3]$
  \item Rate of convergence to normality as a function of the tail index
  \item Edgeworth expansions for finite-sample corrections
  \item Normal as $\alpha = 2$ member of the stable family
\end{itemize}

%#############################################################################
\section{Section 2.4: Stylized Facts of Financial Returns}
%#############################################################################

\subsection{Year 3 Core}
\begin{enumerate}
  \item \textbf{Fat tails / excess kurtosis}: $\kappa = \E[(X-\mu)^4]/\sigma^4 > 3$; typical daily equity kurtosis: $[5, 50]$
  \item \textbf{Negative skewness}: Left tail heavier than right, especially for equity indices
  \item \textbf{Volatility clustering}: $\Corr(|r_t|, |r_{t+k}|) > 0$ for large $k$; ``large changes followed by large changes'' (Mandelbrot, 1963)
  \item \textbf{Aggregational Gaussianity}: Kurtosis decreases as holding period increases (daily $\to$ weekly $\to$ monthly)
\end{enumerate}

\subsection{Master Extension}
\begin{enumerate}[resume]
  \item \textbf{Leverage effect}: Negative returns increase future volatility more than positive returns of equal magnitude
  \item \textbf{Absence of return autocorrelation}: $\Corr(r_t, r_{t+k}) \approx 0$ for $k \geq 1$, yet $\Corr(r_t^2, r_{t+k}^2) > 0$
  \item \textbf{Long memory in volatility}: Slow decay of $\Corr(|r_t|, |r_{t+k}|)$ — hyperbolic, not exponential
  \item \textbf{Gain/loss asymmetry}: Crashes are sharper than rallies
\end{enumerate}

\subsection{PhD Extension}
\begin{itemize}
  \item Cont (2001) complete taxonomy of 11 stylized facts
  \item Multifractal models (Mandelbrot, Fisher \& Calvet, 1997)
  \item Multiscaling: $\E[|r_t|^q] \propto \Delta t^{\zeta(q)}$ with nonlinear $\zeta(q)$
  \item Cross-correlation structure and multivariate stylized facts
  \item Intraday stylized facts: microstructure noise, bid-ask bounce, Epps effect
\end{itemize}

%#############################################################################
\section{Section 2.5: The Student-$t$ Distribution}
%#############################################################################

\subsection{Year 3 Core}

\begin{defn}[Student-$t$ Distribution]
$X \sim t_\nu$ if:
\[
  f(x;\nu) = \frac{\Gamma\!\left(\frac{\nu+1}{2}\right)}{\sqrt{\nu\pi}\,\Gamma\!\left(\frac{\nu}{2}\right)} \left(1 + \frac{x^2}{\nu}\right)^{-(\nu+1)/2}, \quad x \in \R.
\]
\end{defn}

\noindent\textbf{Key properties}:
\begin{itemize}
  \item Single parameter $\nu > 0$ (degrees of freedom) controls tail heaviness
  \item $\E[X] = 0$ ($\nu > 1$); \quad $\Var(X) = \frac{\nu}{\nu-2}$ ($\nu > 2$)
  \item Excess kurtosis $= \frac{6}{\nu - 4}$ ($\nu > 4$)
  \item $\nu \to \infty$: converges to $N(0,1)$; \quad $\nu = 1$: Cauchy distribution
  \item Tails: $f(x) \sim |x|^{-(\nu+1)}$ (power law)
  \item Financial returns: typically $\hat{\nu} \in [3, 8]$ for daily equity data
\end{itemize}

\subsection{Master Extension}
\begin{itemize}
  \item \textbf{Location-scale form}: $X = \mu + \sigma T$ where $T \sim t_\nu$
  \item MLE estimation of $(\mu, \sigma, \nu)$
  \item Information criteria: AIC $= -2\ell + 2k$, BIC $= -2\ell + k\ln n$
  \item Student-$t$ as innovation distribution in GARCH models
  \item QQ-plot diagnostics: systematic deviations indicate misspecification
\end{itemize}

\subsection{PhD Extension}
\begin{itemize}
  \item \textbf{Mixture representation}: $X \mid V \sim N(0, V)$, $V \sim \text{Inv-Gamma}(\nu/2, \nu/2)$
  \item Bayesian interpretation: Student-$t$ as marginal posterior under conjugate normal-inverse-gamma prior
  \item Multivariate Student-$t$: PDF involves $\Gamma$ functions and determinant
  \item Comparison with stable distributions: both have power-law tails, but Student-$t$ has finite moments for $\nu > p$
  \item Tail equivalence: $t_\nu$ tail index $= \nu$ vs.\ stable tail index $= \alpha$
\end{itemize}

%#############################################################################
\section{Section 2.6: Asymmetric Distributions}
%#############################################################################

\subsection{Year 3 Core}
\begin{itemize}
  \item Negative skewness in equity returns: downside risk is larger than upside potential
  \item \textbf{Skewness} measures asymmetry: $S = \E[(X-\mu)^3] / \sigma^3$
  \item Symmetric distributions (Normal, Student-$t$) cannot distinguish left from right tail risk
  \item Intuition: we need distributions with different left and right tail behavior
\end{itemize}

\subsection{Master Extension}

\begin{defn}[Hansen's Skewed-$t$ Distribution (Hansen, 1994)]
\[
  f(x;\nu,\lambda) = \begin{cases}
    bc\!\left(1 + \frac{1}{\nu-2}\left(\frac{bx+a}{1-\lambda}\right)^{\!2}\right)^{-(\nu+1)/2} & x < -a/b \\[6pt]
    bc\!\left(1 + \frac{1}{\nu-2}\left(\frac{bx+a}{1+\lambda}\right)^{\!2}\right)^{-(\nu+1)/2} & x \geq -a/b
  \end{cases}
\]
where $\lambda \in (-1,1)$ is the skewness parameter, $a = 4\lambda c\frac{\nu-2}{\nu-1}$, $b^2 = 1 + 3\lambda^2 - a^2$, $c = \frac{\Gamma((\nu+1)/2)}{\sqrt{\pi(\nu-2)}\,\Gamma(\nu/2)}$.
\end{defn}

\begin{itemize}
  \item $\lambda = 0$: symmetric Student-$t$; \quad $\lambda < 0$: heavier left tail
  \item \textbf{GED}: $f(x;\nu) = \frac{\nu}{2\Gamma(1/\nu)}\exp(-|x|^\nu)$; \quad $\nu = 2$: Normal; $\nu = 1$: Laplace
  \item Widely used as GARCH innovation distributions
\end{itemize}

\subsection{PhD Extension}
\begin{itemize}
  \item \textbf{Generalized Hyperbolic (GH)} distribution: five parameters $(\lambda, \alpha, \beta, \delta, \mu)$
  \item Special cases: Hyperbolic ($\lambda=1$), NIG ($\lambda=-1/2$)
  \item GH as normal mean-variance mixture: $X \mid W \sim N(\mu + \beta W, W)$, $W \sim \text{GIG}(\lambda, \delta, \gamma)$
  \item Barndorff-Nielsen (1977): empirical fit to financial data
  \item Affine closure: GH distributions are closed under affine transformations
\end{itemize}

%#############################################################################
\section{Section 2.7: Stable ($\alpha$-Stable) Distributions \textbf{[MANDATORY]}}
%#############################################################################

\subsection{Year 3 Core}

\begin{defn}[Stable Distribution — Intuitive]
A distribution is \textbf{stable} if a sum of i.i.d.\ copies has the same shape (up to location and scale). Formally, $X \sim S(\alpha, \beta, \gamma, \delta)$ if for i.i.d.\ copies $X_1, X_2$:
\[
  aX_1 + bX_2 \stackrel{d}{=} cX + d, \quad c = (a^\alpha + b^\alpha)^{1/\alpha}.
\]
\end{defn}

\begin{itemize}
  \item \textbf{Four parameters}:
  \begin{itemize}
    \item $\alpha \in (0, 2]$: \textbf{stability index} (tail heaviness) — smaller $\alpha$ $\Rightarrow$ heavier tails
    \item $\beta \in [-1, 1]$: \textbf{skewness}
    \item $\gamma > 0$: \textbf{scale}
    \item $\delta \in \R$: \textbf{location}
  \end{itemize}
  \item \textbf{Special cases}: $\alpha = 2$ (Gaussian), $\alpha = 1, \beta = 0$ (Cauchy)
  \item \textbf{Why variance may not exist}: For $\alpha < 2$, $\E[X^2] = \infty$. The sample variance does not converge.
  \item \textbf{Financial motivation}: Mandelbrot (1963) observed that cotton price changes follow stable laws with $\alpha \approx 1.7$
  \item Typical estimates for financial returns: $\hat{\alpha} \in [1.5, 1.9]$
\end{itemize}

\subsection{Master Extension}

\noindent\textbf{Characteristic function} ($\alpha \neq 1$):
\[
  \phi(t) = \E[e^{itX}] = \exp\!\left\{-\gamma^\alpha |t|^\alpha\!\left(1 - i\beta\,\text{sign}(t)\tan\frac{\pi\alpha}{2}\right) + i\delta t\right\}
\]
\textbf{Case $\alpha = 1$}:
\[
  \phi(t) = \exp\!\left\{-\gamma|t|\!\left(1 + i\beta\frac{2}{\pi}\,\text{sign}(t)\ln|t|\right) + i\delta t\right\}
\]

\begin{itemize}
  \item \textbf{No closed-form PDF} in general; density via inverse Fourier transform:
  \[
    f(x) = \frac{1}{2\pi}\int_{-\infty}^{\infty}\phi(t)\,e^{-itx}\,dt
  \]
  \item \textbf{Generalized Central Limit Theorem (GCLT)}: Stable distributions are the \textit{only} possible non-degenerate limits of normalized sums $\frac{S_n - b_n}{a_n} \xrightarrow{d} S(\alpha, \beta, \gamma, \delta)$
  \item \textbf{Stability property}: Closed under convolution — sums of i.i.d.\ stable r.v.'s remain stable with the same $\alpha$
  \item \textbf{Infinite variance} when $\alpha < 2$: $\E[|X|^p] < \infty$ iff $p < \alpha$
  \item \textbf{Power-law tails}: $P(X > x) \sim C_\alpha(1+\beta)\,x^{-\alpha}$ as $x \to \infty$
  \item \textbf{Estimation methods}: McCulloch quantile method, maximum likelihood (Nolan, 1997), characteristic function methods (Kogon \& Williams, 1998)
  \item \textbf{Financial interpretation}: $\alpha$ measures tail heaviness; $\alpha \approx 1.7$ for S\&P 500 daily returns
\end{itemize}

\subsection{PhD Extension}

\begin{defn}[Regular Variation]
A measurable function $L : (0,\infty) \to (0,\infty)$ is \textbf{slowly varying} at infinity if $\lim_{x\to\infty} L(tx)/L(x) = 1$ for all $t > 0$. A function $f$ is \textbf{regularly varying} with index $\rho$ if $f(x) = x^\rho L(x)$.
\end{defn}

\begin{thm}[Domain of Attraction]
A distribution $F$ belongs to the \textbf{domain of attraction} of a stable law $S(\alpha, \beta, 1, 0)$ if and only if:
\begin{enumerate}
  \item $\bar{F}(x) = x^{-\alpha}L(x)$ for a slowly varying function $L$ (regularly varying tails)
  \item The balance condition: $\lim_{x\to\infty} \frac{\bar{F}(x)}{\bar{F}(x) + F(-x)} = \frac{1+\beta}{2}$
\end{enumerate}
\end{thm}

\begin{itemize}
  \item \textbf{Lévy processes}: A Lévy process $\{L_t\}$ has stationary independent increments; $L_1 \sim S(\alpha, \beta, \gamma, \delta)$ defines a stable Lévy process
  \item \textbf{Lévy--Itô decomposition}: $L_t = bt + \sigma W_t + \int_{|x|\leq 1} x\,\tilde{N}(t,dx) + \int_{|x|>1} x\,N(t,dx)$
  \item \textbf{Lévy measure} of a stable process: $\nu(dx) = \begin{cases} c_+ x^{-\alpha-1}dx & x > 0 \\ c_- |x|^{-\alpha-1}dx & x < 0 \end{cases}$
  \item \textbf{Estimation challenges}:
  \begin{itemize}
    \item MLE requires numerical evaluation of the density (Nolan's algorithm)
    \item Likelihood surface can be multimodal for small samples
    \item Tail index estimators: Hill, Pickands, moment estimator
  \end{itemize}
  \item \textbf{Tempered stable distributions}: Modify the Lévy measure to ensure finite moments: $\nu(dx) = c\,|x|^{-\alpha-1}e^{-\lambda|x|}dx$; retain heavy-tail behavior for moderate $x$ but guarantee finite variance
  \item \textbf{Stable vs.\ Student-$t$}: Both have power-law tails, but Student-$t$ has finite moments up to order $\lfloor\nu\rfloor - 1$; stable has no finite moments beyond order $\alpha$. Student-$t$ is not closed under convolution; stable is. Debate: Mandelbrot--Fama (stable) vs.\ Praetz--Blattberg--Gonedes (Student-$t$)
\end{itemize}

%#############################################################################
\section{Section 2.8: Extreme Value Theory (EVT)}
%#############################################################################

\subsection{Year 3 Core}
\begin{itemize}
  \item \textbf{Motivation}: We need accurate models for extreme events (financial crises, market crashes)
  \item \textbf{Key idea}: Instead of modeling the \textit{entire} distribution, EVT focuses on the \textbf{tails} only
  \item \textbf{GPD concept} (descriptive): Exceedances above a high threshold tend to follow a specific distribution shape — the Generalized Pareto Distribution
  \item \textbf{Tail index $\xi$} (intuitive): a single number that summarizes how ``heavy'' the tail is
  \begin{itemize}
    \item $\xi > 0$: heavy tails (power-law decay) — typical for financial data
    \item $\xi = 0$: exponential tails (similar to Normal)
    \item $\xi < 0$: bounded tails (finite maximum)
  \end{itemize}
  \item Financial data: $\hat{\xi} > 0$ confirms that extreme events are more frequent than the Normal predicts
  \item \textbf{Note}: The formal theorems and VaR computation from EVT are presented at the Master level
\end{itemize}

\subsection{Master Extension}

\begin{thm}[Fisher--Tippett--Gnedenko]
If $M_n = \max(X_1, \ldots, X_n)$ and $\frac{M_n - b_n}{a_n} \xrightarrow{d} H$, then $H$ is the GEV:
\[
  H(x;\xi,\mu,\sigma) = \exp\!\left\{-\left(1 + \xi\frac{x-\mu}{\sigma}\right)^{-1/\xi}\right\}
\]
with $\xi = 0$ (Gumbel), $\xi > 0$ (Fréchet), $\xi < 0$ (Weibull).
\end{thm}

\begin{thm}[Pickands--Balkema--de Haan]
For high threshold $u$, exceedances $Y = X - u \mid X > u$ follow the GPD:
\[
  G(y;\xi,\sigma) = 1 - \left(1 + \xi\frac{y}{\sigma}\right)^{-1/\xi}, \quad y > 0.
\]
\end{thm}

\begin{itemize}
  \item \textbf{POT method}: Estimate $(\xi, \sigma)$ from exceedances above threshold $u$
  \item \textbf{VaR from GPD}:
  \[
    \text{VaR}_\alpha = u + \frac{\sigma}{\xi}\left[\left(\frac{n}{N_u}\cdot\alpha\right)^{-\xi} - 1\right]
  \]
  where $N_u$ = number of exceedances, $n$ = total observations
  \item \textbf{ES from GPD}: $\text{ES}_\alpha = \frac{\text{VaR}_\alpha}{1-\xi} + \frac{\sigma - \xi u}{1-\xi}$, for $\xi < 1$
  \item \textbf{Threshold selection}: Mean excess plot, parameter stability plot
  \item \textbf{Mean excess function}: $e(u) = \E[X - u \mid X > u]$; linear in $u$ for GPD
\end{itemize}

\subsection{PhD Extension}

\begin{defn}[Hill Estimator]
For order statistics $X_{(1)} \geq X_{(2)} \geq \cdots \geq X_{(n)}$, the Hill estimator of the tail index is:
\[
  \hat{\xi}_{\text{Hill}}^{(k)} = \frac{1}{k}\sum_{i=1}^{k}\ln X_{(i)} - \ln X_{(k+1)}, \quad k = 1, \ldots, n-1.
\]
The tail index is $\alpha = 1/\hat{\xi}_{\text{Hill}}$.
\end{defn}

\begin{itemize}
  \item \textbf{Hill plot}: $\hat{\xi}_{\text{Hill}}^{(k)}$ vs.\ $k$; stable region indicates appropriate $k$
  \item \textbf{Bias--variance trade-off}: Small $k$ $\Rightarrow$ high variance; large $k$ $\Rightarrow$ high bias
  \item \textbf{Second-order conditions}: Rate of convergence of the tail to the Pareto form; Hall (1982) second-order parameter $\rho$ governs optimal $k$
  \item \textbf{Adaptive threshold selection}: Minimizing asymptotic MSE (Danielsson et al., 2001)
  \item \textbf{Pickands estimator}: $\hat{\xi} = \frac{1}{\ln 2}\ln\frac{X_{(k)} - X_{(2k)}}{X_{(2k)} - X_{(4k)}}$ — does not require $\xi > 0$
  \item \textbf{Moment estimator} (Dekkers--Einmahl--de Haan): Works for all $\xi \in \R$
  \item \textbf{Connection to stable distributions}: If $X \sim S(\alpha, \beta, \gamma, \delta)$ with $\alpha < 2$, then $\xi = 1/\alpha$
\end{itemize}

%#############################################################################
\section{Section 2.9: Distribution Comparison}
%#############################################################################

\subsection{Year 3 Core}
\begin{itemize}
  \item Histogram with fitted overlays (Normal, Student-$t$, Stable)
  \item QQ-plots: departures from the 45$^\circ$ line
  \item Log-log tail plot: empirical survival function vs.\ theoretical
  \item Visual assessment: which distribution best captures the center, shoulders, and tails?
\end{itemize}

\subsection{Master Extension}
\begin{itemize}
  \item AIC/BIC comparison table across distributions
  \item Kolmogorov--Smirnov test: $D_n = \sup_x |F_n(x) - F_\theta(x)|$
  \item Anderson--Darling test: weighted CDF distance, more sensitive to tails
  \item Likelihood ratio test for nested models (Normal vs.\ Student-$t$)
\end{itemize}

\subsection{PhD Extension}
\begin{itemize}
  \item Vuong (1989) test for non-nested model comparison
  \item Scoring rules: CRPS (Continuous Ranked Probability Score)
  \item Diebold--Mariano test for comparing predictive densities
  \item Cross-validation for distributional models
  \item Probability Integral Transform (PIT) and Rosenblatt test
\end{itemize}

%#############################################################################
\section{Section 2.10: Risk Management Implications}
%#############################################################################

\subsection{Year 3 Core — Introducing VaR from Scratch}

\noindent\textbf{Pedagogical note}: Students at Year 3 level have \textbf{not} encountered VaR or ES before. This section introduces these concepts from first principles.

\vspace{0.3cm}

\noindent\textbf{Step 1: The Risk Question.}\\
A bank holds a portfolio worth \$1{,}000{,}000. The fundamental risk question is:
\begin{quote}
\textit{``What is the maximum loss we can expect on a given day, with a specified level of confidence?''}
\end{quote}

\noindent\textbf{Step 2: Quantile Intuition.}\\
If we know the distribution of daily returns $r_t$, then the $\alpha$-quantile $q_\alpha$ is the value such that:
\[
  P(r_t \leq q_\alpha) = \alpha.
\]
For example, if $\alpha = 0.01$, then $q_{0.01}$ is the return level that is exceeded (on the downside) only 1\% of the time. In 100 trading days, we expect to lose \textit{more} than $q_{0.01}$ on approximately 1 day.

\vspace{0.2cm}

\noindent\textbf{Step 3: Formal Definition.}

\begin{defn}[Value-at-Risk — First Introduction]
The \textbf{Value-at-Risk} at confidence level $(1-\alpha)$ is the loss threshold that is exceeded with probability $\alpha$:
\[
  \text{VaR}_\alpha = -q_\alpha = -F^{-1}(\alpha),
\]
where $F$ is the CDF of returns. The minus sign converts the negative quantile into a positive loss amount.
\end{defn}

\noindent\textbf{Interpretation}: ``With probability $(1-\alpha)$, the daily loss will \textbf{not exceed} $\text{VaR}_\alpha$.''

\vspace{0.2cm}

\noindent\textbf{Step 4: VaR under the Normal Assumption.}\\
If $r_t \sim N(\mu, \sigma^2)$, then:
\[
  \text{VaR}_\alpha = -(\mu + \sigma \cdot z_\alpha), \quad \text{where } z_\alpha = \Phi^{-1}(\alpha).
\]
\textbf{Example}: $\mu = 0$, $\sigma = 0.015$ (daily), portfolio = \$1M, $\alpha = 0.01$:
\[
  \text{VaR}_{0.01} = -0.015 \times (-2.326) \times 1{,}000{,}000 = \$34{,}890.
\]
Interpretation: ``On 99\% of trading days, the loss will not exceed \$34{,}890.''

\vspace{0.2cm}

\noindent\textbf{Step 5: Why the Distribution Matters.}\\
Under a Student-$t$ with $\nu = 5$, the 1\% quantile is more extreme than $-2.326$. So VaR is \textbf{higher}:

\begin{center}
\renewcommand{\arraystretch}{1.2}
\begin{tabular}{l c c c}
\toprule
\textbf{Model} & \textbf{99\% VaR (\$1M)} & \textbf{99.9\% VaR (\$1M)} & \textbf{Ratio to Normal} \\
\midrule
Normal & \$27{,}000 & \$35{,}900 & 1.0$\times$ \\
Student-$t$ ($\nu=5$) & \$33{,}700 & \$58{,}400 & 1.6$\times$ \\
Stable ($\alpha=1.7$) & \$41{,}000 & \$90{,}200 & 2.5$\times$ \\
\bottomrule
\end{tabular}
\end{center}

\begin{itemize}
  \item At the 99\% level, the difference is moderate ($\sim$25\%)
  \item At the 99.9\% level, the difference is \textbf{dramatic}: the Normal underestimates risk by a factor of 1.6--2.5$\times$
  \item \textbf{Conclusion}: The choice of distribution directly affects how much capital a bank must hold
\end{itemize}

\noindent\textbf{Step 6: Graphical Intuition.}\\
A slide should show the left tail of the PDF with a shaded area = $\alpha$, and the VaR marked as the boundary point. Two overlaid distributions (Normal vs.\ Student-$t$) make the difference visually obvious.

\vspace{0.3cm}

\noindent\textbf{Step 7: Limitations of VaR (preview).}\\
\begin{itemize}
  \item VaR tells us the \textit{threshold} of the worst $\alpha$-fraction of days, but not \textit{how bad} things get beyond that threshold
  \item VaR does not distinguish between ``losing \$35k'' and ``losing \$200k'' — both are simply ``VaR exceedances''
  \item A better measure (Expected Shortfall) will be introduced at the Master level
\end{itemize}

\subsection{Master Extension — Expected Shortfall and Backtesting}

\begin{defn}[Expected Shortfall — First Introduction at Master Level]
The \textbf{Expected Shortfall} (also called Conditional VaR) at level $\alpha$ is the average loss in the worst $\alpha$-fraction of scenarios:
\[
  \text{ES}_\alpha = -\frac{1}{\alpha}\int_0^\alpha F^{-1}(p)\,dp = \E[-X \mid X \leq -\text{VaR}_\alpha].
\]
\end{defn}

\noindent\textbf{Interpretation}: ``Given that we are in the worst $\alpha\%$ of days, what is the \textit{average} loss?''

\begin{itemize}
  \item ES is a \textbf{coherent} risk measure (subadditive): $\text{ES}(X+Y) \leq \text{ES}(X) + \text{ES}(Y)$
  \item VaR is \textbf{not} subadditive — diversification can increase VaR
  \item Basel III/IV: shift from 99\% VaR to 97.5\% ES
  \item Kupiec test: $LR_{uc} = -2\ln\frac{\alpha^x(1-\alpha)^{n-x}}{\hat{\alpha}^x(1-\hat{\alpha})^{n-x}} \sim \chi^2(1)$
  \item Christoffersen test: independence of exceedances
  \item Basel traffic light system: Green (0--4), Yellow (5--9), Red (10+) exceedances per 250 days
\end{itemize}

\subsection{PhD Extension}
\begin{itemize}
  \item \textbf{Spectral risk measures}: $\rho_\phi(X) = -\int_0^1 \phi(p)\,F^{-1}(p)\,dp$ with risk spectrum $\phi$
  \item \textbf{Elicitability}: VaR is elicitable (can be backtested via scoring functions); ES is not elicitable alone but jointly elicitable with VaR (Fissler \& Ziegel, 2016)
  \item \textbf{Capital allocation}: Euler allocation under ES — distributional assumptions directly affect allocated capital per desk
  \item \textbf{Stable distributions and extreme risk}: $\text{ES}_\alpha$ under stable distributions diverges as $\alpha \to 0$ faster than under Student-$t$; implications for capital adequacy
  \item \textbf{Model risk quantification}: Range of risk measures across plausible distributional models
\end{itemize}

%#############################################################################
\section{Section 2.11: Modern Extensions}
%#############################################################################

\subsection{Year 3 Core}
\begin{itemize}
  \item GARCH intuition: variance changes over time; yesterday's large shock $\Rightarrow$ today's high variance
  \item $\sigma_t^2 = \omega + \alpha\epsilon_{t-1}^2 + \beta\sigma_{t-1}^2$ (conceptual formula)
  \item After GARCH filtering, standardized residuals are closer to Gaussian but still show excess kurtosis
\end{itemize}

\subsection{Master Extension}
\begin{itemize}
  \item GARCH(1,1) with Student-$t$ innovations: $r_t = \mu + \sigma_t z_t$, $z_t \sim t_\nu$
  \item Gaussian mixtures: $f(x) = \sum_{k=1}^K \pi_k\,\phi(x;\mu_k, \sigma_k^2)$; calm vs.\ crisis regimes
  \item EM algorithm for mixture estimation
  \item Hamilton (1989) regime-switching: Markov chain latent state $s_t$
  \item Asymmetric GARCH: GJR-GARCH, EGARCH for leverage effect
\end{itemize}

\subsection{PhD Extension}
\begin{itemize}
  \item Stochastic volatility: $d\ln\sigma_t^2 = \kappa(\theta - \ln\sigma_t^2)\,dt + \sigma_v\,dW_t^v$ with correlated Brownian motions
  \item Jump-diffusion models: Merton (1976), Kou (2002)
  \item Lévy-driven volatility models: BNS model (Barndorff-Nielsen \& Shephard, 2001)
  \item Realized volatility and high-frequency estimation
  \item GARCH with stable innovations: Mittnik \& Rachev (1993)
  \item Fractionally integrated GARCH (FIGARCH) for long memory
\end{itemize}


%=============================================================================
%=============================================================================
\part{KEY FORMULAS}
%=============================================================================
%=============================================================================

\section{Complete Formula Reference}

\subsection{Returns and Moments}

\begin{enumerate}
  \item \textbf{Log-return}: $r_t = \ln(P_t / P_{t-1})$

  \item \textbf{Sample mean}: $\bar{r} = \frac{1}{n}\sum_{t=1}^n r_t$

  \item \textbf{Sample variance}: $\hat{\sigma}^2 = \frac{1}{n-1}\sum_{t=1}^n (r_t - \bar{r})^2$

  \item \textbf{Skewness}: $S = \frac{1}{n}\sum_{t=1}^n \left(\frac{r_t - \bar{r}}{\hat{\sigma}}\right)^3$

  \item \textbf{Kurtosis}: $K = \frac{1}{n}\sum_{t=1}^n \left(\frac{r_t - \bar{r}}{\hat{\sigma}}\right)^4$

  \item \textbf{Excess kurtosis}: $K_e = K - 3$
\end{enumerate}

\subsection{Distribution PDFs}

\begin{enumerate}[resume]
  \item \textbf{Normal PDF}: $f(x) = \frac{1}{\sigma\sqrt{2\pi}}\exp\!\left(-\frac{(x-\mu)^2}{2\sigma^2}\right)$

  \item \textbf{Student-$t$ PDF}: $f(x;\nu) = \frac{\Gamma((\nu+1)/2)}{\sqrt{\nu\pi}\,\Gamma(\nu/2)}\left(1 + \frac{x^2}{\nu}\right)^{-(\nu+1)/2}$

  \item \textbf{Student-$t$ excess kurtosis}: $\kappa_e = \frac{6}{\nu - 4}$, \quad $\nu > 4$

  \item \textbf{Hansen's Skewed-$t$ PDF}: see Section 2.6

  \item \textbf{GED PDF}: $f(x;\nu) = \frac{\nu}{2\Gamma(1/\nu)}\exp(-|x|^\nu)$
\end{enumerate}

\subsection{Stable Distributions}

\begin{enumerate}[resume]
  \item \textbf{Stable characteristic function} ($\alpha \neq 1$):
  \[
    \phi(t) = \exp\!\left\{-\gamma^\alpha|t|^\alpha\!\left(1 - i\beta\,\text{sign}(t)\tan\frac{\pi\alpha}{2}\right) + i\delta t\right\}
  \]

  \item \textbf{Moment existence}: $\E[|X|^p] < \infty$ iff $p < \alpha$

  \item \textbf{Tail behavior}: $P(X > x) \sim C_\alpha(1+\beta)\,x^{-\alpha}$, \quad $C_\alpha = \frac{1-\alpha}{\Gamma(2-\alpha)\cos(\pi\alpha/2)}$
\end{enumerate}

\subsection{Extreme Value Theory}

\begin{enumerate}[resume]
  \item \textbf{GEV CDF}: $H(x;\xi,\mu,\sigma) = \exp\!\left\{-\left(1 + \xi\frac{x-\mu}{\sigma}\right)^{-1/\xi}\right\}$

  \item \textbf{GPD CDF}: $G(y;\xi,\sigma) = 1 - \left(1 + \xi\frac{y}{\sigma}\right)^{-1/\xi}$, \quad $y > 0$

  \item \textbf{Hill estimator}: $\hat{\xi}_{\text{Hill}}^{(k)} = \frac{1}{k}\sum_{i=1}^{k}\ln X_{(i)} - \ln X_{(k+1)}$

  \item \textbf{Mean excess function}: $e(u) = \E[X - u \mid X > u] = \frac{\sigma + \xi u}{1 - \xi}$ (GPD)
\end{enumerate}

\subsection{Risk Measures}

\begin{enumerate}[resume]
  \item \textbf{VaR} \layerY{(introduced in Section 2.10)}: $\text{VaR}_\alpha = -F^{-1}(\alpha)$

  \item \textbf{VaR under normality} \layerY{}: $\text{VaR}_\alpha = -(\mu + \sigma\,\Phi^{-1}(\alpha))$

  \item \textbf{ES under normality} \layerM{}: $\text{ES}_\alpha = -\mu + \sigma\frac{\phi(\Phi^{-1}(\alpha))}{\alpha}$

  \item \textbf{VaR from GPD} \layerM{}: $\text{VaR}_\alpha = u + \frac{\sigma}{\xi}\left[\left(\frac{n\alpha}{N_u}\right)^{-\xi} - 1\right]$

  \item \textbf{ES from GPD} \layerM{}: $\text{ES}_\alpha = \frac{\text{VaR}_\alpha}{1-\xi} + \frac{\sigma - \xi u}{1-\xi}$

  \item \textbf{Kupiec LR test} \layerM{}: $LR_{uc} = -2\ln\frac{\alpha^x(1-\alpha)^{n-x}}{\hat{\alpha}^x(1-\hat{\alpha})^{n-x}} \sim \chi^2(1)$
\end{enumerate}

\subsection{Model Selection}

\begin{enumerate}[resume]
  \item \textbf{AIC}: $\text{AIC} = -2\ell + 2k$

  \item \textbf{BIC}: $\text{BIC} = -2\ell + k\ln n$

  \item \textbf{Jarque--Bera}: $JB = \frac{n}{6}\left(S^2 + \frac{(K-3)^2}{4}\right) \sim \chi^2(2)$
\end{enumerate}

All symbols are defined where first introduced:
\begin{itemize}[nosep]
  \item $P_t$: asset price; $r_t$: log-return; $n$: sample size
  \item $\mu$: location/mean; $\sigma$: scale/std dev; $\nu$: degrees of freedom
  \item $\alpha$: stability index (stable) or tail probability (VaR)
  \item $\beta$: skewness (stable); $\gamma$: scale (stable); $\delta$: location (stable)
  \item $\xi$: shape parameter (GEV/GPD); $\ell$: log-likelihood; $k$: number of parameters
  \item $\Phi$: standard normal CDF; $\phi$: standard normal PDF; $\Gamma$: gamma function
\end{itemize}


%=============================================================================
%=============================================================================
\part{EMPIRICAL ILLUSTRATIONS}
%=============================================================================
%=============================================================================

\section{Required Graphs — Mapping to Existing Charts}

\begin{center}
\renewcommand{\arraystretch}{1.4}
\begin{tabular}{c l l l}
\toprule
\textbf{\#} & \textbf{Required Graph} & \textbf{Existing File} & \textbf{Layer} \\
\midrule
G1 & Price series + log-return plot & \texttt{ch2\_normal\_fit.png} & Y3 \\
G2 & Histogram + Normal fit overlay & \texttt{ch2\_normal\_fit.png} & Y3 \\
G3 & Histogram + Student-$t$ fit & \texttt{ch2\_student\_t\_fit.png} & Y3 \\
G4 & Student-$t$ PDFs (varying $\nu$) & \texttt{ch2\_student\_t\_pdfs.png} & Y3 \\
G5 & Histogram + Stable fit comparison & \texttt{ch2\_sp500\_stable\_fit.png} & Y3/MSc \\
G6 & Stable PDFs (varying $\alpha$) & \texttt{ch2\_stable\_pdfs.png} & MSc \\
G7 & QQ-plot vs.\ Normal & \texttt{ch2\_normality\_tests.png} & Y3 \\
G8 & QQ-plot vs.\ Student-$t$ & \texttt{ch2\_dist\_comparison\_qq.png} & Y3 \\
G9 & Log-log tail plot & \texttt{ch2\_dist\_comparison\_tails.png} & MSc \\
G10 & Distribution overlay comparison & \texttt{ch2\_dist\_comparison\_overlay.png} & Y3 \\
G11 & CLT illustration & \texttt{ch2\_clt\_illustration.png} & Y3 \\
G12 & GEV fit to block maxima & \texttt{ch2\_evt\_gev\_fit.png} & MSc \\
G13 & GPD / POT analysis & \texttt{ch2\_evt\_gpd\_pot.png} & MSc \\
G14 & EVT-based VaR and ES & \texttt{ch2\_evt\_var\_es.png} & MSc \\
G15 & VaR comparison bar chart & \texttt{ch2\_var\_comparison.png} & Y3 \\
G16 & VaR backtest & \texttt{ch2\_var\_backtest.png} & MSc \\
G17 & Risk comparison table & \texttt{ch2\_risk\_comparison\_table.png} & Y3/MSc \\
G18 & GARCH-$t$ fit & \texttt{ch2\_garch\_t\_fit.png} & MSc \\
G19 & Cross-asset stability analysis & \texttt{ch2\_cross\_asset\_alpha.png} & MSc/PhD \\
G20 & Tail survival probability plot & \texttt{ch2\_tail\_survival.png} & PhD \\
\bottomrule
\end{tabular}
\end{center}

\subsection{Additional Graphs Needed}

\begin{center}
\renewcommand{\arraystretch}{1.3}
\begin{tabular}{c l l}
\toprule
\textbf{\#} & \textbf{Graph} & \textbf{Layer} \\
\midrule
G21 & Hill plot ($\hat{\xi}$ vs.\ $k$) & PhD \\
G22 & Mean excess plot with threshold line & MSc/PhD \\
G23 & Stable vs.\ Student-$t$ tail comparison (log-log) & PhD \\
G24 & VaR graphical intuition: PDF with shaded tail and VaR threshold (TikZ) & Y3 \\
G24b & VaR exceedance count by model (bar chart) & Y3/MSc \\
G25 & Aggregational Gaussianity: kurtosis vs.\ horizon & Y3/MSc \\
G26 & QQ-plot vs.\ Stable distribution & MSc \\
G27 & Regime classification plot & MSc \\
G28 & Gaussian mixture fit & MSc \\
\bottomrule
\end{tabular}
\end{center}


%=============================================================================
%=============================================================================
\part{QUANTLET PLAN}
%=============================================================================
%=============================================================================

\section{Quantlet Architecture}

\subsection{Q1: Stylized Facts Analysis}

\begin{tabular}{l p{12cm}}
\toprule
\textbf{Name} & \texttt{SFM\_ch2\_stylized\_facts} \\
\textbf{Layer} & \layerY{Year 3 Core} \\
\textbf{Objective} & Demonstrate the key stylized facts of financial returns: fat tails, skewness, volatility clustering, aggregational Gaussianity \\
\textbf{Data} & S\&P 500 daily prices from Yahoo Finance (\texttt{yfinance}), 2000--2025 \\
\textbf{Steps} &
  1. Download prices, compute log-returns \\
  & 2. Plot price series and return series side by side \\
  & 3. Compute sample moments: mean, std, skewness, kurtosis \\
  & 4. Histogram with Normal overlay \\
  & 5. QQ-plot vs.\ Normal \\
  & 6. ACF of returns vs.\ ACF of $|r_t|$ and $r_t^2$ (volatility clustering) \\
  & 7. Kurtosis at different horizons (daily, weekly, monthly) \\
\textbf{Graphs} & G1, G2, G7, G11, G25 \\
\textbf{Interpretation} & Report: normality rejected (Jarque--Bera), kurtosis $\gg 3$, volatility clustering evident, convergence to normality at lower frequencies \\
\bottomrule
\end{tabular}

\subsection{Q2: Normal vs.\ Student-$t$ Fit}

\begin{tabular}{l p{12cm}}
\toprule
\textbf{Name} & \texttt{SFM\_ch2\_normal\_vs\_student} \\
\textbf{Layer} & \layerY{Year 3 Core} \\
\textbf{Objective} & Fit Normal and Student-$t$ distributions via MLE; compare goodness of fit \\
\textbf{Data} & S\&P 500 daily log-returns \\
\textbf{Steps} &
  1. MLE fit of $N(\mu, \sigma^2)$ \\
  & 2. MLE fit of location-scale $t_\nu(\mu, \sigma)$ \\
  & 3. Report parameters: $(\hat{\mu}, \hat{\sigma})$ and $(\hat{\mu}, \hat{\sigma}, \hat{\nu})$ \\
  & 4. Compute AIC, BIC for both models \\
  & 5. Histogram with both fitted PDFs overlaid \\
  & 6. QQ-plots: Normal and Student-$t$ \\
\textbf{Graphs} & G2, G3, G4, G7, G8, G10 \\
\textbf{Interpretation} & Student-$t$ achieves substantially better AIC/BIC despite extra parameter; QQ-plot improvement in tails; typical $\hat{\nu} \in [3, 6]$ \\
\bottomrule
\end{tabular}

\subsection{Q3: VaR Comparison}

\begin{tabular}{l p{12cm}}
\toprule
\textbf{Name} & \texttt{SFM\_ch2\_var\_comparison} \\
\textbf{Layer} & \layerY{Year 3} / \layerM{Master} \\
\textbf{Objective} & Compute and compare VaR and ES under Normal, Student-$t$, and empirical quantile approaches \\
\textbf{Data} & S\&P 500 daily log-returns \\
\textbf{Steps} &
  1. Fit Normal and Student-$t$ distributions \\
  & 2. Compute VaR at 95\%, 99\%, 99.9\% under each model \\
  & 3. Compute empirical VaR (historical simulation) \\
  & 4. Compute ES under each model \\
  & 5. Bar chart comparing all VaR estimates \\
  & 6. Table with \$1M portfolio loss estimates \\
  & 7. [MSc] Rolling window VaR backtest (250-day window) \\
  & 8. [MSc] Kupiec test for each model \\
\textbf{Graphs} & G15, G16, G17, G24 \\
\textbf{Interpretation} & Normal severely underestimates tail risk at extreme quantiles; Student-$t$ and empirical approaches are more conservative; backtesting confirms Normal has too many violations \\
\bottomrule
\end{tabular}

\subsection{Q4: Tail Estimation (Hill Estimator)}

\begin{tabular}{l p{12cm}}
\toprule
\textbf{Name} & \texttt{SFM\_ch2\_tail\_estimation} \\
\textbf{Layer} & \layerM{Master} / \layerP{PhD} \\
\textbf{Objective} & Estimate the tail index using the Hill estimator and compare with parametric tail indices \\
\textbf{Data} & S\&P 500 daily absolute log-returns \\
\textbf{Steps} &
  1. Sort absolute returns in descending order \\
  & 2. Compute Hill estimator $\hat{\xi}^{(k)}$ for $k = 10, 11, \ldots, 500$ \\
  & 3. Plot Hill plot: $\hat{\alpha} = 1/\hat{\xi}$ vs.\ $k$ \\
  & 4. Identify stable region for tail index estimate \\
  & 5. Compare with $\hat{\alpha}$ from stable fit and $\hat{\nu}$ from Student-$t$ fit \\
  & 6. [PhD] Pickands and moment estimators for comparison \\
  & 7. [PhD] Bias-corrected Hill estimator \\
\textbf{Graphs} & G20, G21, G23 \\
\textbf{Interpretation} & Hill estimate typically gives $\hat{\alpha} \in [3, 5]$ (corresponding to Student-$t$ with $\nu \approx 3$--$5$); stable fit gives $\hat{\alpha} \in [1.5, 1.9]$; discuss the difference between tail index and stability index \\
\bottomrule
\end{tabular}

\subsection{Q5: Stable Distribution Simulation and Fitting}

\begin{tabular}{l p{12cm}}
\toprule
\textbf{Name} & \texttt{SFM\_ch2\_stable\_distributions} \\
\textbf{Layer} & \layerY{Year 3} / \layerM{Master} \\
\textbf{Objective} & Simulate from stable distributions; fit stable model to S\&P 500 returns \\
\textbf{Data} & Simulated data + S\&P 500 daily returns \\
\textbf{Steps} &
  1. Simulate $S(\alpha, 0, 1, 0)$ for $\alpha = 1.2, 1.5, 1.7, 2.0$ \\
  & 2. Plot PDFs on linear and log scales \\
  & 3. Demonstrate stability property: $X_1 + X_2 \stackrel{d}{=} cX + d$ \\
  & 4. [MSc] Fit stable distribution to S\&P 500 via MLE (\texttt{scipy.stats.levy\_stable}) \\
  & 5. [MSc] Report $(\hat{\alpha}, \hat{\beta}, \hat{\gamma}, \hat{\delta})$ \\
  & 6. [MSc] Compare stable fit vs.\ Normal on log-scale plot \\
  & 7. [MSc] Cross-asset analysis: estimate $\hat{\alpha}$ for multiple assets \\
\textbf{Graphs} & G5, G6, G19, G26 \\
\textbf{Interpretation} & Stable distributions capture heavy tails better than Normal; estimated $\hat{\alpha} < 2$ confirms infinite-variance regime; discuss implications for portfolio theory \\
\bottomrule
\end{tabular}

\subsection{Q6: EVT-Based VaR}

\begin{tabular}{l p{12cm}}
\toprule
\textbf{Name} & \texttt{SFM\_ch2\_evt\_var} \\
\textbf{Layer} & \layerM{Master} / \layerP{PhD} \\
\textbf{Objective} & Apply EVT (POT method) to estimate tail risk measures \\
\textbf{Data} & S\&P 500 daily losses (negated returns) \\
\textbf{Steps} &
  1. Compute losses $L_t = -r_t$ \\
  & 2. Mean excess plot: $\hat{e}(u)$ vs.\ $u$ \\
  & 3. Select threshold $u$ from linearity region \\
  & 4. Fit GPD to exceedances via MLE \\
  & 5. Report $(\hat{\xi}, \hat{\sigma})$; confirm $\hat{\xi} > 0$ \\
  & 6. Compute EVT-based VaR and ES at 99\%, 99.9\% \\
  & 7. Compare with Normal, Student-$t$, and historical simulation \\
  & 8. [PhD] GEV fit to monthly block maxima \\
  & 9. [PhD] Hill plot and tail index comparison \\
  & 10. [PhD] Parameter stability plot across thresholds \\
\textbf{Graphs} & G12, G13, G14, G22 \\
\textbf{Interpretation} & EVT provides theoretically justified extrapolation beyond observed data; GPD-based VaR is more conservative than Normal but less extreme than Stable; discuss threshold sensitivity \\
\bottomrule
\end{tabular}


%=============================================================================
%=============================================================================
\part{EXERCISES}
%=============================================================================
%=============================================================================

\section{Exercises — Three Difficulty Levels}

%#############################################################################
\subsection{Year 3 Exercises (Undergraduate)}
%#############################################################################

\begin{enumerate}[label=\textbf{E\arabic*.}]

\item \textbf{Log-Returns Computation.}
Download 5 years of daily closing prices for Apple (AAPL) from Yahoo Finance.
\begin{enumerate}[label=(\alph*)]
  \item Compute the daily log-returns.
  \item Plot the price series and the return series.
  \item Compute the sample mean, standard deviation, skewness, and kurtosis.
  \item Is the excess kurtosis significantly different from zero? What does this imply?
\end{enumerate}

\item \textbf{Normal vs.\ Student-$t$ Fit.}
Using the AAPL returns from Exercise 1:
\begin{enumerate}[label=(\alph*)]
  \item Fit a Normal distribution by MLE and report $(\hat{\mu}, \hat{\sigma})$.
  \item Fit a Student-$t$ distribution by MLE and report $(\hat{\mu}, \hat{\sigma}, \hat{\nu})$.
  \item Produce a histogram with both fitted PDFs overlaid.
  \item Produce QQ-plots against both distributions. Which fits better in the tails?
  \item Compute AIC and BIC for both models. Which model is preferred?
\end{enumerate}

\item \textbf{Introduction to VaR.}
This exercise introduces Value-at-Risk (VaR) step by step.
\begin{enumerate}[label=(\alph*)]
  \item In your own words, explain what the statement ``the 99\% VaR is \$25{,}000'' means for a bank.
  \item Using the AAPL returns from Exercise 1, compute the sample mean $\hat{\mu}$ and standard deviation $\hat{\sigma}$.
  \item Under the Normal assumption, the 1\% quantile is $q_{0.01} = \hat{\mu} + \hat{\sigma} \times (-2.326)$. Compute this value.
  \item The VaR is $\text{VaR}_{0.01} = -q_{0.01}$. Compute it. For a \$100{,}000 portfolio, what is the maximum daily loss at 99\% confidence?
  \item Now compute the 1\% \textit{empirical} quantile directly from the data (i.e., sort all returns and find the value below which 1\% of observations fall). Compare with the Normal-based VaR. Which is larger? What does the difference tell you about the Normal assumption?
\end{enumerate}

\item \textbf{Stylized Facts.}
Using S\&P 500 daily returns (2000--2024):
\begin{enumerate}[label=(\alph*)]
  \item Compute the kurtosis at daily, weekly, and monthly frequencies.
  \item Does the kurtosis decrease with the holding period? This is called \emph{aggregational Gaussianity} — explain why it occurs.
  \item Plot the ACF of returns and the ACF of absolute returns. What is the key difference? Relate to volatility clustering.
\end{enumerate}

\item \textbf{Stable Distributions: Conceptual.}
\begin{enumerate}[label=(\alph*)]
  \item What does it mean for a distribution to be ``stable''?
  \item Name three special cases of the stable distribution family.
  \item If $\alpha = 1.7$, does the variance exist? Does the mean exist? Explain.
  \item Why did Mandelbrot (1963) advocate stable distributions for financial returns?
\end{enumerate}

\end{enumerate}

%#############################################################################
\subsection{Master Exercises}
%#############################################################################

\begin{enumerate}[label=\textbf{M\arabic*.}]

\item \textbf{Comprehensive Distribution Fitting.}
Download 20 years of daily returns for a major equity index (e.g., S\&P 500, DAX, or Nikkei 225).
\begin{enumerate}[label=(\alph*)]
  \item Fit the following distributions by MLE: Normal, Student-$t$, Skewed-$t$ (Hansen), Stable.
  \item Report all fitted parameters in a table.
  \item Compute AIC, BIC, and log-likelihood for each model.
  \item Produce a panel of QQ-plots (one per distribution).
  \item On a log-scale histogram, overlay all fitted PDFs. Which model captures the tails best?
  \item Discuss: Is the Stable distribution's infinite variance a practical problem?
\end{enumerate}

\item \textbf{VaR Backtesting.}
Using a rolling window of 250 trading days:
\begin{enumerate}[label=(\alph*)]
  \item Compute the 1-day ahead 99\% VaR under Normal, Student-$t$, and historical simulation.
  \item Count the number of VaR exceedances for each model over the out-of-sample period.
  \item Apply the Kupiec test to each model. Which models pass at the 5\% significance level?
  \item Apply the Christoffersen test. Are exceedances independent?
  \item Classify each model under the Basel traffic light system.
\end{enumerate}

\item \textbf{EVT Application.}
\begin{enumerate}[label=(\alph*)]
  \item For S\&P 500 daily losses, plot the mean excess function.
  \item Select an appropriate threshold and justify your choice.
  \item Fit the GPD to exceedances and report $(\hat{\xi}, \hat{\sigma})$.
  \item Compute the EVT-based 99\% and 99.9\% VaR and ES.
  \item Compare with Normal-based and Student-$t$-based estimates.
  \item Fit the GEV distribution to monthly block maxima and compare $\hat{\xi}$.
\end{enumerate}

\item \textbf{GARCH-$t$ Filtering.}
\begin{enumerate}[label=(\alph*)]
  \item Fit a GARCH(1,1) model with Normal innovations to S\&P 500 returns.
  \item Fit a GARCH(1,1) model with Student-$t$ innovations.
  \item Compare the models using AIC/BIC and the log-likelihood.
  \item Extract standardized residuals from the GARCH-$t$ model. Are they closer to Gaussian than the raw returns?
  \item Compute the kurtosis of the standardized residuals. How much of the original excess kurtosis is ``explained'' by GARCH?
\end{enumerate}

\item \textbf{Cross-Asset Analysis.}
Estimate the stability index $\hat{\alpha}$ for five different assets (e.g., S\&P 500, EUR/USD, Gold, Bitcoin, 10-Year Treasury).
\begin{enumerate}[label=(\alph*)]
  \item Report the estimated parameters in a table.
  \item Which asset class has the heaviest tails?
  \item Is $\hat{\alpha}$ significantly different from 2 for any of the assets?
  \item Discuss implications for portfolio diversification.
\end{enumerate}

\end{enumerate}

%#############################################################################
\subsection{PhD Exercises}
%#############################################################################

\begin{enumerate}[label=\textbf{P\arabic*.}]

\item \textbf{Tail Index Estimation.}
For S\&P 500 daily absolute returns:
\begin{enumerate}[label=(\alph*)]
  \item Compute the Hill estimator $\hat{\xi}^{(k)}$ for $k = 10, \ldots, 500$.
  \item Plot the Hill plot and identify a stable region.
  \item Compute the Pickands and moment estimators. Do they agree?
  \item Derive the asymptotic variance of the Hill estimator: $\Var(\hat{\xi}^{(k)}) \sim \xi^2/k$ under second-order conditions.
  \item Apply the Danielsson et al.\ (2001) adaptive method for optimal $k$ selection.
  \item Compare $\hat{\alpha}_{\text{Hill}} = 1/\hat{\xi}$ with the $\hat{\alpha}$ from the stable fit. Discuss.
\end{enumerate}

\item \textbf{Domain of Attraction.}
\begin{enumerate}[label=(\alph*)]
  \item Show that the Student-$t_\nu$ distribution is in the domain of attraction of a stable law with index $\alpha = \nu$.
  \item Prove that the normal distribution is in the domain of attraction of $S(2, 0, \gamma, 0)$.
  \item For the Cauchy distribution, verify directly that $\bar{F}(x) \sim (\pi x)^{-1}$ is regularly varying with index $-1$.
  \item Discuss: If returns follow a Student-$t_5$, what is the appropriate normalizing sequence $(a_n, b_n)$ for the maximum?
\end{enumerate}

\item \textbf{Lévy Process Simulation.}
\begin{enumerate}[label=(\alph*)]
  \item Simulate a symmetric stable Lévy process $L_t$ with $\alpha = 1.7$ at times $t = 1/252, 2/252, \ldots, 1$.
  \item Verify the self-similarity property: $L_{ct} \stackrel{d}{=} c^{1/\alpha} L_t$.
  \item Simulate a tempered stable process with exponential tempering $e^{-\lambda|x|}$ for $\lambda = 0.1, 1, 10$.
  \item Compare the tail behavior of the tempered stable with the pure stable.
  \item Compute sample moments at different time scales; verify that the tempered stable has finite variance but the pure stable does not.
\end{enumerate}

\item \textbf{Scoring Rules and Model Comparison.}
\begin{enumerate}[label=(\alph*)]
  \item Implement the Continuous Ranked Probability Score (CRPS) for each fitted distribution.
  \item Compute the log-score (predictive log-likelihood) out of sample.
  \item Apply the Diebold--Mariano test to compare Normal vs.\ Student-$t$ predictive densities.
  \item Apply the Vuong test to compare Student-$t$ vs.\ Stable (non-nested models).
  \item Implement the Probability Integral Transform (PIT) and test for uniformity.
\end{enumerate}

\item \textbf{Spectral Risk Measures.}
\begin{enumerate}[label=(\alph*)]
  \item Define the spectral risk measure $\rho_\phi(X) = -\int_0^1 \phi(p)F^{-1}(p)\,dp$ for a risk spectrum $\phi$.
  \item For the exponential spectrum $\phi(p) = \frac{ke^{-k(1-p)}}{1-e^{-k}}$, compute $\rho_\phi$ under Normal and Student-$t$ assumptions.
  \item Show that VaR and ES are special cases of spectral risk measures.
  \item Discuss elicitability: why is ES not elicitable alone? What does joint elicitability (Fissler \& Ziegel, 2016) mean for backtesting?
\end{enumerate}

\end{enumerate}


%=============================================================================
%=============================================================================
\part{EXAM QUESTIONS}
%=============================================================================
%=============================================================================

\section{Year 3 Exam Questions}

\begin{enumerate}[label=\textbf{Y3-Q\arabic*.}]

\item \textbf{(25 points)} Consider daily log-returns of a stock with sample statistics: $\bar{r} = 0.0003$, $\hat{\sigma} = 0.015$, skewness $= -0.4$, kurtosis $= 7.2$.

\begin{enumerate}[label=(\alph*)]
  \item (5 pts) Define the log-return and explain why it is preferred over simple returns in statistical analysis.
  \item (5 pts) Compute the 99\% VaR under the Normal assumption for a \$1{,}000{,}000 portfolio. Use $z_{0.01} = -2.326$.
  \item (5 pts) The excess kurtosis is $7.2 - 3 = 4.2$. What does this imply about the tails of the return distribution?
  \item (5 pts) A Student-$t$ distribution is fitted with $\hat{\nu} = 5$. Verify that the theoretical excess kurtosis is $6/(5-4) = 6$. Why might the sample excess kurtosis (4.2) differ from this?
  \item (5 pts) Which distributional model (Normal or Student-$t$) would you recommend for computing risk measures? Justify.
\end{enumerate}

\item \textbf{(25 points)}

\begin{enumerate}[label=(\alph*)]
  \item (8 pts) State the Central Limit Theorem precisely. What conditions are required?
  \item (5 pts) Explain ``aggregational Gaussianity'' in the context of financial returns.
  \item (7 pts) A Student-$t$ distribution with $\nu = 3$ has infinite kurtosis. Does the CLT apply to sums of i.i.d.\ $t_3$ random variables? Explain carefully. (Hint: the CLT requires finite variance.)
  \item (5 pts) What is the Generalized CLT, and how does it relate to stable distributions?
\end{enumerate}

\item \textbf{(25 points)}

\begin{enumerate}[label=(\alph*)]
  \item (7 pts) Define Value-at-Risk at the $\alpha$ level. What does the statement ``the 99\% VaR is \$50{,}000'' mean in plain language? Draw a diagram showing a PDF with the VaR threshold and the shaded tail area.
  \item (8 pts) For a portfolio with returns $\sim N(0.001, 0.02^2)$, compute the 95\% and 99\% VaR for a \$1M investment. Use $z_{0.05} = -1.645$, $z_{0.01} = -2.326$. Show all steps.
  \item (5 pts) Now suppose the returns follow a Student-$t$ with $\nu = 5$, same location and scale. Without computing exact values, explain \textit{qualitatively} why the VaR would be higher than under the Normal assumption. (Hint: compare the tails of the two distributions.)
  \item (5 pts) State one important limitation of VaR as a risk measure. (Hint: VaR tells you the threshold of the worst $\alpha\%$ of days, but what does it \textit{not} tell you?)
\end{enumerate}

\end{enumerate}

\section{Master Exam Questions}

\begin{enumerate}[label=\textbf{MSc-Q\arabic*.}]

\item \textbf{(30 points)} You are fitting distributions to 5{,}000 daily returns of the DAX index.

\begin{enumerate}[label=(\alph*)]
  \item (5 pts) Write the log-likelihood function for the location-scale Student-$t$ model with parameters $(\mu, \sigma, \nu)$.
  \item (5 pts) The Normal fit gives $\ell_N = 14{,}200$ (with $k = 2$). The Student-$t$ fit gives $\ell_t = 14{,}450$ (with $k = 3$). Compute the AIC and BIC for both models ($n = 5{,}000$). Which model is preferred?
  \item (8 pts) Write the characteristic function of the stable distribution $S(\alpha, \beta, \gamma, \delta)$ for $\alpha \neq 1$. What is the stability property? Prove that $S + S' \stackrel{d}{=} c S + d$ for i.i.d.\ stable copies using the characteristic function.
  \item (7 pts) Explain the Peaks Over Threshold (POT) method. Write the GPD CDF and derive the VaR formula from the GPD.
  \item (5 pts) You estimate $\hat{\xi} = 0.25$ from the GPD fit. What does this imply about the tail behavior? Is the variance of the exceedances finite?
\end{enumerate}

\item \textbf{(30 points)} VaR backtesting and regulatory framework.

\begin{enumerate}[label=(\alph*)]
  \item (5 pts) Define the Kupiec (unconditional coverage) test. State the null hypothesis, the test statistic, and its asymptotic distribution.
  \item (8 pts) Over 250 trading days, a bank's 99\% VaR model produces 8 exceedances. Compute the Kupiec test statistic. Is the model rejected at the 5\% level? ($\chi^2_{0.05}(1) = 3.84$.)
  \item (7 pts) Classify this result under the Basel traffic light system. What are the regulatory consequences?
  \item (5 pts) What additional information does the Christoffersen (independence) test provide beyond the Kupiec test?
  \item (5 pts) Explain why Expected Shortfall was adopted over VaR in Basel III/IV. What is ``coherence'' and why does VaR fail it?
\end{enumerate}

\item \textbf{(30 points)} Distribution comparison and risk management.

\begin{enumerate}[label=(\alph*)]
  \item (8 pts) You estimate the following for S\&P 500 returns: $\hat{\alpha}_{\text{stable}} = 1.72$, $\hat{\nu}_{\text{Student-}t} = 4.8$, $\hat{\xi}_{\text{GPD}} = 0.22$. Explain what each parameter measures. Are the three estimates of ``tail heaviness'' consistent? (Hint: for the GPD, $\xi \approx 1/\alpha_{\text{tail}}$.)
  \item (7 pts) Compute the 99.9\% VaR under the Normal assumption with $\mu = 0$, $\sigma = 0.012$, using $z_{0.001} = -3.09$. Then compute it under the Student-$t$ with $\nu = 5$ using $t_{5, 0.001} = -4.77$. Compute the ratio.
  \item (8 pts) Explain the difference between the tail index estimated by the Hill estimator and the stability index $\alpha$ of a stable distribution. Can both be estimated from the same data? What do they measure differently?
  \item (7 pts) A GARCH(1,1) model with Normal innovations produces standardized residuals with excess kurtosis 1.5 (down from 18 in the raw returns). Interpret this result. Does it justify using Student-$t$ innovations?
\end{enumerate}

\end{enumerate}

\section{PhD Exam Questions}

\begin{enumerate}[label=\textbf{PhD-Q\arabic*.}]

\item \textbf{(40 points)} Stable distributions and domains of attraction.

\begin{enumerate}[label=(\alph*)]
  \item (8 pts) State the Generalized Central Limit Theorem. Under what conditions does the limit of normalized sums converge to a stable distribution with $\alpha < 2$?
  \item (10 pts) Prove that if $F$ has a regularly varying tail with index $-\alpha$ (i.e., $\bar{F}(x) = x^{-\alpha}L(x)$), then $F$ is in the domain of attraction of a stable law $S(\alpha, \beta, 1, 0)$.
  \item (10 pts) Consider the Student-$t$ distribution with $\nu$ degrees of freedom. Show that $P(X > x) \sim c_\nu\,x^{-\nu}$ as $x \to \infty$ and identify the constant $c_\nu$. Hence determine the domain of attraction of $t_\nu$.
  \item (7 pts) Explain the Mandelbrot--Fama vs.\ Praetz--Blattberg--Gonedes debate. What are the practical implications of choosing stable vs.\ Student-$t$ for risk management?
  \item (5 pts) Define tempered stable distributions. How do they resolve the infinite-variance issue while retaining heavy-tail behavior?
\end{enumerate}

\item \textbf{(40 points)} Extreme value theory and tail estimation.

\begin{enumerate}[label=(\alph*)]
  \item (8 pts) State the Fisher--Tippett--Gnedenko theorem. Derive the three types (Gumbel, Fréchet, Weibull) as special cases of the GEV.
  \item (10 pts) Prove the Pickands--Balkema--de Haan theorem: for a distribution $F$ in the domain of attraction of $H_\xi$, the excess distribution $F_u(y) = P(X - u \leq y \mid X > u)$ converges to $G_{\xi, \sigma(u)}(y)$ as $u \to x^*$.
  \item (8 pts) Derive the asymptotic distribution of the Hill estimator: $\sqrt{k}(\hat{\xi}_H^{(k)} - \xi) \xrightarrow{d} N(0, \xi^2)$ under appropriate conditions. State the required second-order regular variation condition.
  \item (7 pts) Explain the bias--variance trade-off in threshold selection for the Hill estimator. How does the Danielsson--de Haan--Peng--Vries (2001) method address this?
  \item (7 pts) Discuss the connection between the GPD shape parameter $\xi$ and the stable tail index $\alpha$: if $X \sim S(\alpha, \beta, \gamma, \delta)$, what is $\xi$ in the GPD approximation to the upper tail?
\end{enumerate}

\item \textbf{(40 points)} Advanced risk measurement.

\begin{enumerate}[label=(\alph*)]
  \item (8 pts) Define a \textbf{coherent risk measure} (Artzner et al., 1999). State the four axioms. Show that VaR violates subadditivity with a counterexample.
  \item (8 pts) Define spectral risk measures. Show that ES is a spectral risk measure with constant spectrum $\phi(p) = 1/\alpha$ for $p \leq \alpha$.
  \item (8 pts) Discuss elicitability in the context of risk measure backtesting. Why is ES not elicitable alone? State the Fissler--Ziegel (2016) joint elicitability result for (VaR, ES).
  \item (8 pts) For a portfolio with returns following a stable distribution $S(1.7, -0.2, 0.01, 0)$, discuss the computation of VaR and ES. What special difficulties arise from the lack of closed-form PDF? How can they be addressed numerically?
  \item (8 pts) Compare the regulatory approaches of Basel II (99\% VaR) and Basel III/IV (97.5\% ES). Under which distributional assumptions do these two approaches yield approximately the same capital requirement? Under which do they differ most?
\end{enumerate}

\end{enumerate}


%=============================================================================
%=============================================================================
\part{RISK MANAGEMENT CONNECTION}
%=============================================================================
%=============================================================================

\section{Underestimation of Tail Risk Under Normality}

\begin{center}
\renewcommand{\arraystretch}{1.3}
\begin{tabular}{l c c c}
\toprule
\textbf{Event} & \textbf{$\sigma$-Multiple} & \textbf{$P$ under Normal} & \textbf{Empirical Frequency} \\
\midrule
3$\sigma$ event & 3.0 & 0.13\% & $\sim$1\% (daily) \\
5$\sigma$ event & 5.0 & $3 \times 10^{-7}$ & $\sim$0.1\% (daily) \\
10$\sigma$ event & 10.0 & $\sim 10^{-23}$ & Observed (Oct 1987) \\
25$\sigma$ event & 25.0 & $\sim 10^{-135}$ & Observed (Oct 1987) \\
\bottomrule
\end{tabular}
\end{center}

\section{Impact on VaR and ES}

\begin{itemize}
  \item At the \textbf{95\% level}: All models give similar VaR estimates; distribution choice matters little.
  \item At the \textbf{99\% level}: Student-$t$ gives 20--30\% higher VaR than Normal.
  \item At the \textbf{99.9\% level}: Stable gives 2.5$\times$ higher VaR; EVT gives 2.1$\times$ higher VaR.
  \item \textbf{ES amplifies the differences}: ES captures the \textit{severity} of tail losses, not just their threshold.
  \item Under stable distributions with $\alpha < 2$, ES can be dramatically larger because the conditional tail expectation is infinite in the limit.
\end{itemize}

\section{Stable Distributions and Extreme Risk}

\begin{itemize}
  \item If returns truly follow a stable distribution with $\alpha < 2$:
  \begin{itemize}
    \item The variance does not exist — traditional Markowitz optimization is meaningless.
    \item The CLT does not apply — aggregation does not lead to normality.
    \item Instead, the GCLT applies: sums converge to a stable law.
    \item Risk measures based on variance (Sharpe ratio, tracking error) are unreliable.
  \end{itemize}
  \item \textbf{Practical resolution}: Tempered stable distributions maintain heavy-tail behavior for moderate extremes while ensuring finite variance.
\end{itemize}

\section{Implications for Capital Allocation}

\begin{itemize}
  \item Under-capitalization under Normal assumptions can lead to systemic fragility.
  \item Basel III/IV shift to ES partially addresses this by focusing on tail severity.
  \item The choice between Student-$t$, Stable, and EVT-based risk measures directly affects:
  \begin{itemize}
    \item Market risk capital requirements
    \item Desk-level risk limits
    \item Stress testing scenarios
    \item Credit risk modeling (default correlation in the tails)
  \end{itemize}
  \item \textbf{Recommendation}: Use Student-$t$ or GARCH-$t$ for day-to-day risk management; use EVT for stress testing and extreme quantile estimation; consider stable distributions for theoretical understanding and research.
\end{itemize}


%=============================================================================
%=============================================================================
\part{REFERENCES}
%=============================================================================
%=============================================================================

\section{Core References}

\begin{enumerate}
  \item Mandelbrot, B. (1963). The Variation of Certain Speculative Prices. \textit{Journal of Business}, 36(4), 394--419.
  \item Fama, E.F. (1965). The Behavior of Stock-Market Prices. \textit{Journal of Business}, 38(1), 34--105.
  \item Cont, R. (2001). Empirical Properties of Asset Returns: Stylized Facts and Statistical Issues. \textit{Quantitative Finance}, 1(2), 223--236.
  \item Franke, J., Härdle, W.K., \& Hafner, C.M. (2019). \textit{Statistics of Financial Markets}. 4th ed. Springer.
  \item McNeil, A.J., Frey, R., \& Embrechts, P. (2015). \textit{Quantitative Risk Management}. Rev.\ ed. Princeton UP.
  \item Nolan, J.P. (2020). \textit{Stable Distributions: Models for Heavy Tailed Data}. Birkhäuser.
  \item Embrechts, P., Klüppelberg, C., \& Mikosch, T. (1997). \textit{Modelling Extremal Events}. Springer.
  \item Bollerslev, T. (1986). Generalized Autoregressive Conditional Heteroskedasticity. \textit{J.\ Econometrics}, 31(3), 307--327.
  \item Hansen, B.E. (1994). Autoregressive Conditional Density Estimation. \textit{Int.\ Econ.\ Rev.}, 35(3), 705--730.
  \item Hamilton, J.D. (1989). A New Approach to Nonstationary Time Series. \textit{Econometrica}, 57(2), 357--384.
  \item Artzner, P., Delbaen, F., Eber, J.M., \& Heath, D. (1999). Coherent Measures of Risk. \textit{Math.\ Finance}, 9(3), 203--228.
  \item Fissler, T. \& Ziegel, J.F. (2016). Higher Order Elicitability and Osband's Principle. \textit{Ann.\ Statist.}, 44(4), 1680--1707.
  \item Danielsson, J., de Haan, L., Peng, L., \& de Vries, C. (2001). Using a Bootstrap Method to Choose the Sample Fraction in Tail Index Estimation. \textit{J.\ Multivariate Analysis}, 76, 226--248.
\end{enumerate}

\end{document}
